\chapter{Inventory Policies}
\label{chap:policies}

Seven inventory control policies are implemented in the \texttt{policies/}
package.  They share a common Python interface and cover the major families
found in the MEIO literature, from continuous-review analytical policies to
periodic-review heuristics and the echelon-optimal Clark-Scarf policy.

%% -----------------------------------------------------------------------
\section{Common API}
\label{sec:policies:api}

Every policy exposes a single method:

\begin{lstlisting}[language=Python, caption={Universal policy interface.}]
def order_qty(
    self,
    on_hand: int,
    backlog_external: int,
    backlog_children: int,
    pipeline_in: int,
    t: Optional[int] = None,
) -> int:
    ...
\end{lstlisting}

The \texttt{t} argument is the current simulation day.  Periodic policies
inspect \texttt{t} to determine whether a review is due; continuous-review
policies ignore it.  Passing \texttt{t=None} causes periodic policies to
act as if a review is due, which is the convention used during warm-up.

The \emph{inventory position} (IP) used in most policies is:

\begin{equation}
  \text{IP} = \text{on\_hand} - \text{backlog\_external}
             + \text{pipeline\_in}
  \label{eq:ip}
\end{equation}

Orders are always non-negative; over-stocking is never triggered.

%% -----------------------------------------------------------------------
\section{Base-Stock Policy}
\label{sec:policies:bs}

\textbf{Implementation:} \texttt{policies/base\_stock.py}

The base-stock (order-up-to) policy is the single-parameter benchmark
derived from newsvendor theory.  On every period, the order quantity is:

\begin{equation}
  q = \max\!\bigl(0,\; S - \text{IP}\bigr)
  \label{eq:basestock}
\end{equation}

where $S$ is the \emph{base-stock level}.  The optimal $S$ under Gaussian
demand with mean $\mu$ and standard deviation $\sigma$ over a lead time $L$
is the newsvendor quantile:

\begin{equation}
  S^* = \mu(L+1) + z_\alpha \,\sigma\sqrt{L+1}
  \label{eq:newsvendor}
\end{equation}

where $z_\alpha$ is the standard-normal quantile corresponding to the
target fill rate $\alpha$ (e.g.\ $z_{0.92} \approx 1.41$).  This formula
is used to initialise the analytical baseline in experiment E0.

The base-stock policy is \emph{optimal for single-echelon, uncapacitated,
backlog models} but does not account for transport capacity, vehicle
consolidation, or order batching.

%% -----------------------------------------------------------------------
\section{Reorder-Point Family: $(s,S)$ and $(R,Q)$}
\label{sec:policies:sS}

\textbf{Implementations:} \texttt{policies/ss\_policy.py},
\texttt{policies/rq\_policy.py}

\subsection{$(s,S)$ policy}

The $(s,S)$ policy is a continuous-review policy with two parameters:
a reorder point $s$ and an order-up-to level $S > s$.  An order is
placed only when the inventory position falls at or below $s$:

\begin{equation}
  q =
  \begin{cases}
    S - \text{IP} & \text{if } \text{IP} \leq s \\
    0             & \text{otherwise}
  \end{cases}
\end{equation}

Unlike the base-stock policy, $(s,S)$ batches demand between review events,
which naturally aligns order sizes with vehicle loads.  This batching
behaviour was one motivation for including it alongside base-stock in the
preliminary validation experiments (Section~\ref{sec:meth:ev3}).

\subsection{$(R,Q)$ policy}

The $(R,Q)$ (reorder-point, fixed-quantity) policy places a fixed
order of size $Q$ whenever the inventory position drops to or below the
reorder point $R$:

\begin{equation}
  q =
  \begin{cases}
    Q & \text{if } \text{IP} \leq R \\
    0 & \text{otherwise}
  \end{cases}
\end{equation}

Because $Q$ is fixed, orders are always the same size — well-suited to
lanes with a fixed vehicle capacity.

%% -----------------------------------------------------------------------
\section{Periodic-Review: $(R,S)$ and Order-Up-To with Phase Offset}
\label{sec:policies:rs}

\textbf{Implementations:} \texttt{policies/order\_up\_to.py},
\texttt{policies/periodic\_review.py}

\subsection{$(R,S)$ periodic-review policy}

The \texttt{PeriodicReviewPolicy} checks inventory every $R$ days and
orders up to level $S$:

\begin{equation}
  q =
  \begin{cases}
    \max(0, S - \text{IP}) & \text{if } t \bmod R = 0 \\
    0                      & \text{otherwise}
  \end{cases}
\end{equation}

Between reviews, no order is placed regardless of the inventory position.
The periodic structure reduces ordering overhead and can be co-ordinated
across nodes to smooth vehicle utilisation.

\subsection{Order-up-to with phase offset}

The \texttt{OrderUpToPolicy} generalises the above by allowing a
\texttt{phase\_offset}: reviews occur when $(t - \text{phase\_offset}) \bmod R = 0$.
This lets different nodes review on staggered days, spreading transport
demand across the week and improving consolidation opportunities.

Optional \texttt{k} and \texttt{m} parameters further restrict the
maximum number of orders ($k$) within a sliding window of $m$ days,
combining periodic review with a $(k,m)$-cycle cap.

%% -----------------------------------------------------------------------
\section{Cyclic Policy: $(k,m)$-Cycle Scheduling}
\label{sec:policies:km}

\textbf{Implementation:} \texttt{policies/km\_cycle.py}

The \texttt{KmCyclePolicy} schedules at most $k$ replenishment orders per
$m$-day cycle, at pre-specified intra-cycle offsets:

\begin{equation}
  q =
  \begin{cases}
    \max(0, S - \text{IP}) & \text{if } t \bmod m \in \texttt{review\_offsets}
                             \text{ and orders this cycle} < k \\
    0                      & \text{otherwise}
  \end{cases}
\end{equation}

For example, with $m = 7$ and \texttt{review\_offsets} $= \{0, 3\}$, the
policy reviews on Monday and Thursday of each week.  This structure
is designed to be \emph{consolidation-friendly}: if the warehouse always
orders on fixed days, the supplier can prepare full truck loads in advance.

%% -----------------------------------------------------------------------
\section{Echelon Base-Stock (Clark--Scarf)}
\label{sec:policies:echelon}

\textbf{Implementation:} \texttt{policies/echelon\_stock.py}

The echelon base-stock policy implements the optimal policy for serial
multi-echelon systems identified by \citet{ClarkScarf1960}.  Rather than
conditioning on the local inventory position, it conditions on the
\emph{echelon inventory position} (EIP) — the stock at the current node
plus all downstream nodes, both on-hand and in-pipeline:

\begin{equation}
  \text{EIP}(n) = \text{on\_hand}(n)
                + \sum_{d \in \text{downstream}(n)}
                  \bigl[\text{on\_hand}(d) + \text{pipeline\_in}(d)\bigr]
  \label{eq:eip}
\end{equation}

The order quantity is:

\begin{equation}
  q = \max\!\bigl(0,\; S_e - \text{EIP}(n)\bigr)
\end{equation}

where $S_e$ is the echelon base-stock level.  The simulator computes
the EIP via a DFS traversal of the downstream subgraph in
\texttt{Simulator.\_echelon\_pipeline}, and passes the result to the
policy through the \texttt{pipeline\_in} argument.

The echelon stock policy is optimal under the restrictive assumptions
of Clark and Scarf (serial network, independent demands, no capacity
constraints), which are \emph{not} fully satisfied in the 1N3 topology.
Nevertheless, it provides an important upper bound on what a classical
analytical approach can achieve.

%% -----------------------------------------------------------------------
\section{Policy Comparison}
\label{sec:policies:comparison}

Table~\ref{tab:policy-summary} summarises the seven policies along four
dimensions: number of parameters, review type, whether the policy
accounts for network structure, and its amenability to mixed-integer
optimisation.

\begin{table}[htbp]
  \centering
  \caption{Summary of the seven implemented inventory policies.}
  \label{tab:policy-summary}
  \renewcommand{\arraystretch}{1.2}
  \begin{tabular}{@{}lcccc@{}}
    \toprule
    Policy & Params & Review & Network-aware & Opt-friendly \\
    \midrule
    Base-stock        & 1 & Continuous & No  & Yes \\
    $(s,S)$           & 2 & Continuous & No  & Yes \\
    $(R,Q)$           & 2 & Continuous & No  & Yes \\
    $(R,S)$ periodic  & 2 & Periodic   & No  & Yes \\
    Order-up-to + offset & 2--4 & Periodic & No & Yes \\
    $(k,m)$-cycle     & 3--4 & Periodic   & No  & Yes \\
    Echelon base-stock & 1 & Continuous & Yes & Yes \\
    \bottomrule
  \end{tabular}
\end{table}

All policies are parameterised by small integer or continuous scalars,
making them amenable to the Optuna-TPE optimiser described in
Chapter~\ref{chap:optimization}.  The echelon base-stock policy is the
only policy that explicitly accounts for downstream inventory, at the cost
of requiring the simulator to compute the EIP on every tick.

The main optimisation experiments (E0--E4) use the base-stock policy
for all nodes because its single-parameter structure keeps the search
space tractable (15 parameters for 5 SKUs $\times$ 3 non-supplier nodes)
while being flexible enough to demonstrate the joint optimisation benefit.
